\begin{table}[H]
\centering
\small
\resizebox{\textwidth}{!}{%
\begin{tabular}{lcccccc}
\toprule
Галактика & $\mu_{150}$ & $D_{150}$, Mpc & $\mu_{\mathrm{TRGB}}$ & $D_{\mathrm{TRGB}}$, Mpc & $\mu_{160,\mathrm{Jensen}}$ & $D_{160,\mathrm{Jensen}}$, Mpc \\
\midrule
NGC~1380 & 31.326 $\pm$ 0.062 & 18.42 $\pm$ 0.53 & 31.408 $\pm$ 0.025 & 19.12 $\pm$ 0.22 & 31.518 $\pm$ 0.128 & 20.12 $\pm$ 1.19 \\
NGC~1399 & 31.535 $\pm$ 0.109 & 20.27 $\pm$ 1.02 & 31.498 $\pm$ 0.026 & 19.93 $\pm$ 0.24 & 31.537 $\pm$ 0.127 & 20.30 $\pm$ 1.19 \\
NGC~1404 & 31.369 $\pm$ 0.063 & 18.78 $\pm$ 0.54 & 31.366 $\pm$ 0.025 & 18.76 $\pm$ 0.22 & 31.488 $\pm$ 0.134 & 19.84 $\pm$ 1.22 \\
NGC~1549 & 31.190 $\pm$ 0.053 & 17.30 $\pm$ 0.42 & 31.231 $\pm$ 0.026 & 17.63 $\pm$ 0.21 & -- & -- \\
NGC~3379 & 30.163 $\pm$ 0.057 & 10.78 $\pm$ 0.28 & 30.078 $\pm$ 0.028 & 10.37 $\pm$ 0.13 & -- & -- \\
NGC~4374 & 31.128 $\pm$ 0.054 & 16.81 $\pm$ 0.42 & 31.112 $\pm$ 0.027 & 16.69 $\pm$ 0.21 & -- & -- \\
NGC~4406 & 31.051 $\pm$ 0.052 & 16.23 $\pm$ 0.39 & 31.096 $\pm$ 0.025 & 16.57 $\pm$ 0.19 & -- & -- \\
NGC~4472 & 30.961 $\pm$ 0.061 & 15.57 $\pm$ 0.43 & 31.016 $\pm$ 0.027 & 15.97 $\pm$ 0.20 & 31.046 $\pm$ 0.132 & 16.19 $\pm$ 0.98 \\
NGC~4486 & 30.980 $\pm$ 0.067 & 15.70 $\pm$ 0.48 & 31.046 $\pm$ 0.025 & 16.19 $\pm$ 0.19 & -- & -- \\
NGC~4552 & 31.005 $\pm$ 0.059 & 15.88 $\pm$ 0.43 & 30.923 $\pm$ 0.028 & 15.30 $\pm$ 0.20 & -- & -- \\
NGC~4621 & 30.906 $\pm$ 0.055 & 15.18 $\pm$ 0.39 & 30.821 $\pm$ 0.026 & 14.59 $\pm$ 0.17 & -- & -- \\
NGC~4636 & 31.107 $\pm$ 0.066 & 16.65 $\pm$ 0.51 & 31.070 $\pm$ 0.026 & 16.37 $\pm$ 0.20 & -- & -- \\
NGC~4649 & 31.110 $\pm$ 0.063 & 16.67 $\pm$ 0.48 & 31.009 $\pm$ 0.025 & 15.91 $\pm$ 0.18 & 31.001 $\pm$ 0.135 & 15.86 $\pm$ 0.98 \\
NGC~4697 & 30.245 $\pm$ 0.065 & 11.19 $\pm$ 0.34 & 30.324 $\pm$ 0.028 & 11.61 $\pm$ 0.15 & -- & -- \\
\bottomrule
\end{tabular}
}
\caption{Сравнение расстояний, полученных по JWST/F150W SBF в данной работе, с литературной JWST/TRGB шкалой и HST/F160W SBF-калибровкой Jensen et al. (2015). Ошибка $\mu_{150}$ включает ошибку измерения $\bar{m}_{150}$ и внутренний разброс цветовой калибровки $\sigma_{\rm int}=0.049$ mag. Прочерки означают отсутствие соответствующего F160W измерения в пересекающейся выборке Jensen et al. (2015).}
\label{tab:distance_comparison}
\end{table}
